\documentclass[12pt,a4paper]{article}
\usepackage[utf8]{inputenc}
\usepackage[T1]{fontenc}
\usepackage{amsmath,amssymb,amsthm}
\usepackage{physics}
\usepackage{geometry}
\usepackage{enumitem}
\usepackage{hyperref}
\usepackage{fancyhdr}
\usepackage{titlesec}

\geometry{margin=1in}
\pagestyle{fancy}
\fancyhf{}
\rhead{Lecture 9}
\lhead{Quantum Mechanics --- The Theoretical Minimum}
\cfoot{\thepage}

\titleformat{\section}{\large\bfseries}{}{0em}{}
\titleformat{\subsection}{\normalsize\bfseries}{}{0em}{}

\newtheorem{definition}{Definition}
\newtheorem{theorem}{Theorem}

\title{\textbf{Lecture 9: Particles, Fourier Analysis, and the Momentum Operator} \\ \large From Spins to Continuous Systems}
\author{Based on Leonard Susskind's \textit{The Theoretical Minimum}}
\date{}

\begin{document}
\maketitle
\thispagestyle{fancy}

% ============================================================
\section{1. Transition to Continuous Systems}
% ============================================================

Up to now we have studied the foundations of quantum mechanics using the simplest discrete system---a spin-$\tfrac{1}{2}$ particle with a two-dimensional state space.  We now move to a particle moving on a one-dimensional line, where the state space is \emph{infinite-dimensional} and the observables (position, momentum) take continuous values.

Classically, a particle's state requires \emph{both} position $x$ and momentum $p$.  In quantum mechanics, specifying both is too much information---$x$ and $p$ do not commute, so you describe the state by one \emph{or} the other.  The transition from discrete to continuous systems is mainly notational: sums become integrals, Kronecker deltas become Dirac deltas, and column vectors become wave functions.

% ============================================================
\section{2. Mathematical Interlude: Fourier Analysis}
% ============================================================

\subsection{2.1 Plane waves}
For a fixed number $p$, the function
\[
    e^{ipx} = \cos(px) + i\sin(px)
\]
is a \textbf{plane wave} with ``wavenumber'' $p$.  Larger $|p|$ means shorter wavelength (faster oscillation).

\subsection{2.2 Fourier transform pair}
A very wide class of well-behaved functions $\psi(x)$ can be decomposed into plane waves:
\[
    \boxed{\psi(x) = \frac{1}{\sqrt{2\pi}}\int_{-\infty}^{\infty} \tilde{\psi}(p)\,e^{ipx}\,dp,}
\]
where $\tilde{\psi}(p)$ is the \textbf{Fourier transform} of $\psi(x)$.  The inverse relation is
\[
    \boxed{\tilde{\psi}(p) = \frac{1}{\sqrt{2\pi}}\int_{-\infty}^{\infty} \psi(x)\,e^{-ipx}\,dx.}
\]
$\tilde{\psi}(p)$ gives the amplitude of the wave component with wavenumber $p$ in the function $\psi(x)$.

\subsection{2.3 Physical interpretation}
In quantum mechanics:
\begin{itemize}[nosep]
    \item $\psi(x)$ is the \textbf{position-space wave function}: $|\psi(x)|^2\,dx$ is the probability of finding the particle between $x$ and $x+dx$.
    \item $\tilde{\psi}(p)$ is the \textbf{momentum-space wave function}: $|\tilde{\psi}(p)|^2\,dp$ is the probability of finding the momentum between $p$ and $p+dp$.
    \item The Fourier transform connects the two representations.
\end{itemize}

% ============================================================
\section{3. The Dirac Delta Function}
% ============================================================

\begin{definition}[Delta function]
The Dirac delta function $\delta(x-x')$ is defined operationally by its action on any reasonable function $f$:
\[
    \int_{-\infty}^{\infty} \delta(x-x')\,f(x)\,dx = f(x').
\]
\end{definition}

\subsection{3.1 Key properties}
\begin{enumerate}[nosep]
    \item $\delta(x-x')\,f(x) = f(x')\,\delta(x-x')$ \quad (only nonzero when $x=x'$).
    \item $\displaystyle\int_{-\infty}^{\infty}\delta(x-x')\,dx = 1$.
    \item Combining 1 and 2: $\displaystyle\int \delta(x-x')\,f(x)\,dx = f(x')$.
\end{enumerate}

\subsection{3.2 Fourier representation of the delta function}
Substitute the Fourier transform $\tilde{\psi}(p)$ into the inverse transform for $\psi(x')$:
\[
    \psi(x') = \frac{1}{\sqrt{2\pi}}\int \tilde{\psi}(p)\,e^{ipx'}dp
    = \frac{1}{\sqrt{2\pi}}\int\left(\frac{1}{\sqrt{2\pi}}\int\psi(x)\,e^{-ipx}dx\right)e^{ipx'}dp.
\]
Interchange the order of integration:
\[
    \psi(x') = \int\psi(x)\underbrace{\left(\frac{1}{2\pi}\int e^{ip(x'-x)}dp\right)}_{\text{must be }\delta(x'-x)}dx.
\]
Comparing with property 3 of the delta function yields the fundamental identity:
\[
    \boxed{\delta(x-x') = \frac{1}{2\pi}\int_{-\infty}^{\infty} e^{ip(x-x')}\,dp.}
\]
By symmetry (interchanging $x\leftrightarrow p$):
\[
    \delta(p-p') = \frac{1}{2\pi}\int_{-\infty}^{\infty} e^{ix(p-p')}\,dx.
\]

% ============================================================
\section{4. Particle on a Line: State Space}
% ============================================================

\subsection{4.1 Position basis}
The position eigenstates $\{\ket{x}\}$ satisfy
\[
    \hat{x}\ket{x} = x\ket{x}, \qquad \braket{x'}{x} = \delta(x-x').
\]
The completeness (resolution of identity) relation is
\[
    \int_{-\infty}^{\infty}\ket{x}\bra{x}\,dx = \mathbf{1}.
\]
Any state can be expanded:
\[
    \ket{\Psi} = \int \psi(x)\,\ket{x}\,dx, \qquad \psi(x) = \braket{x}{\Psi}.
\]

\noindent\textbf{Position operator in wave-function language:} The operator $\hat{x}$ acts on wave functions by multiplication: $\hat{x}\,\psi(x)=x\,\psi(x)$.  The eigenfunctions are delta functions: if the particle is at $x_0$, its wave function is $\psi(x)=\delta(x-x_0)$, and $\hat{x}\,\delta(x-x_0)=x\,\delta(x-x_0)=x_0\,\delta(x-x_0)$, confirming eigenvalue $x_0$.

\subsection{4.2 Momentum basis}
Similarly, momentum eigenstates $\{\ket{p}\}$ satisfy
\[
    \hat{p}\ket{p} = p\ket{p}, \qquad \braket{p'}{p} = \delta(p-p').
\]
The overlap between position and momentum eigenstates is
\[
    \braket{x}{p} = \frac{1}{\sqrt{2\pi}}\,e^{ipx},
\]
which is precisely the plane wave.  The Fourier transform pair follows:
\[
    \psi(x) = \braket{x}{\Psi} = \frac{1}{\sqrt{2\pi}}\int \tilde{\psi}(p)\,e^{ipx}\,dp, \qquad
    \tilde{\psi}(p) = \braket{p}{\Psi} = \frac{1}{\sqrt{2\pi}}\int \psi(x)\,e^{-ipx}\,dx.
\]

% ============================================================
\section{5. The Momentum Operator in Position Space}
% ============================================================

Acting on a position-space wave function, the momentum operator is represented as a derivative:
\[
    \boxed{\hat{p}\,\psi(x) = -i\frac{d}{dx}\psi(x).}
\]

\subsection{5.1 Verification}
A plane wave $\psi(x)=e^{ipx}$ is a momentum eigenstate:
\[
    \hat{p}\,e^{ipx} = -i\frac{d}{dx}e^{ipx} = -i(ip)\,e^{ipx} = p\,e^{ipx}. \quad\checkmark
\]

\subsection{5.2 Canonical commutation relation}
Compute $[\hat{x},\hat{p}]$ acting on an arbitrary $\psi(x)$:
\begin{align}
    \hat{x}\hat{p}\,\psi &= x\bigl(-i\psi'(x)\bigr) = -ix\psi'(x), \\
    \hat{p}\hat{x}\,\psi &= -i\frac{d}{dx}\bigl(x\psi(x)\bigr) = -i\bigl(\psi(x)+x\psi'(x)\bigr).
\end{align}
Subtracting: $(\hat{x}\hat{p}-\hat{p}\hat{x})\psi = -ix\psi'+i\psi+ix\psi' = i\psi$.  Since this holds for all $\psi$:
\[
    \boxed{[\hat{x},\hat{p}] = i\mathbf{1}} \qquad (\text{with } \hbar=1).
\]
This is the fundamental commutation relation of quantum mechanics.  Since $[\hat{x},\hat{p}]\neq 0$, position and momentum \textbf{cannot be simultaneously measured}---the basis of the Heisenberg uncertainty principle.

\noindent\textbf{Restoring $\hbar$:} $[\hat{x},\hat{p}]=i\hbar$.  The commutator is the quantum analogue of the classical Poisson bracket $\{x,p\}=1$, via $\{\cdot,\cdot\}\leftrightarrow\frac{1}{i\hbar}[\cdot,\cdot]$.

% ============================================================
\section{6. Simple Hamiltonians for Particles}
% ============================================================

\subsection{6.1 Relativistic particle: $H=cP$}
The Schr\"odinger equation $i\frac{\partial\psi}{\partial t}=c(-i\frac{\partial\psi}{\partial x})$ simplifies to:
\[
    \frac{\partial\psi}{\partial t} + c\frac{\partial\psi}{\partial x} = 0.
\]
General solution: $\psi(x,t)=f(x-ct)$ for any shape $f$.  The entire probability distribution moves rigidly to the right at speed $c$ without changing shape.  This describes a massless particle (e.g.\ a neutrino or phonon) moving in one direction.

\subsection{6.2 Non-relativistic free particle: $H=\hat{p}^2/2m$}
The standard Schr\"odinger equation for a free particle of mass $m$:
\[
    \boxed{i\hbar\frac{\partial\psi}{\partial t} = -\frac{\hbar^2}{2m}\frac{\partial^2\psi}{\partial x^2}.}
\]
Unlike $H=cP$, waves of different wavelength (momentum) now move at different speeds---the wave function \textbf{disperses} (spreads out) over time.  The expectation value of position still obeys the classical equation $\frac{d}{dt}\expval{x}=\frac{\expval{p}}{m}$, but the wave packet broadens.

% ============================================================
\section{7. Entanglement of Electrons and Photons}
% ============================================================

A brief example: two electrons in a superposition of singlet and triplet states can emit a photon and become entangled.

\begin{itemize}[nosep]
    \item Start with an unentangled state $\ket{\uparrow\downarrow}\otimes\ket{0_\gamma}$ (no photon).
    \item The triplet component decays to the singlet by emitting a photon of energy $\omega$.
    \item The final state: $\ket{\text{singlet}}\otimes\bigl(c_0\ket{0_\gamma}+c_1\ket{1_\gamma}\bigr)$.  The electrons are entangled with each other; the radiation field is in a superposition of zero and one photon.
    \item This illustrates that \textbf{superpositions of different particle numbers} are perfectly valid quantum states.
\end{itemize}

% ============================================================
\section{8. Summary}
% ============================================================

\begin{enumerate}[nosep]
    \item Moving from discrete to continuous systems: sums $\to$ integrals, $\delta_{ij}\to\delta(x-x')$, columns $\to$ wave functions.
    \item The \textbf{Fourier transform} connects position-space and momentum-space wave functions.
    \item The \textbf{Dirac delta function} $\delta(x-x')=\frac{1}{2\pi}\int e^{ip(x-x')}dp$ is the continuous orthogonality relation.
    \item Position eigenstates: $\braket{x'}{x}=\delta(x-x')$; eigenfunctions are delta functions.
    \item Momentum eigenstates: $\braket{x}{p}=\frac{1}{\sqrt{2\pi}}e^{ipx}$ (plane waves, completely delocalized).
    \item The momentum operator is $\hat{p}=-i\frac{d}{dx}$; $[\hat{x},\hat{p}]=i$.
    \item $H=cP$: rigid wave-packet motion at speed $c$ (massless particles).
    \item $H=\hat{p}^2/2m$: the standard Schr\"odinger equation; wave packets disperse.
\end{enumerate}

\end{document}
